\section{报告说明}

本报告说明「寻迹」校园失物招领平台五人团队在需求、设计、实现、测试、演化和交付阶段的分工与贡献。团队采用模块主责、相邻协作和统一集成的工作方式，使前端、后端、AI、工程化与课程文档围绕同一业务闭环协同推进。

成员姓名与 GitHub 身份由团队于 2026 年 7 月 12 日确认；除项目负责人外，报告使用成员姓名和 GitHub 身份进行识别。

\subsection{团队目标与协作原则}

团队以「形成可演示、可测试、可部署的校园失物招领闭环」为共同目标，采用以下原则：

\begin{enumerate}
  \item 按前端、后端 A、后端 B、AI 与工程化文档划分主责，同时要求关键接口由相邻模块成员共同联调；
  \item 需求、枚举、接口和状态机以仓库文档为共享契约，代码变更同步维护文档；
  \item 使用 \term{main/develop/feature/fix} 分支、Pull Request 和持续集成完成协作；
  \item 由项目负责人统一集成，模块负责人对自己负责的功能、测试和问题修复承担闭环责任；
  \item 文档共创、评审与验收由参与成员共同署名，并在最终提交中明确记录归属。
\end{enumerate}

\section{成员与角色}

\begin{longtable}{L{2.3cm}L{3.5cm}L{3.5cm}L{5.1cm}}
\caption{团队成员与主要职责}\label{tab:members}\\
\toprule
成员 & GitHub 身份 & 主要角色 & 责任范围 \\
\midrule
\endfirsthead
\toprule
成员 & GitHub 身份 & 主要角色 & 责任范围 \\
\midrule
\endhead
朱梓涵 & \term{handsomezhuzhu} & 项目负责人、后端 A、集成与 QA & 项目统筹；用户、认证、物品模块；核心闭环修复；CI、部署集成；与程文韬共同完成并提交文档 \\
周星辰 & \term{stevezxc} & 前端负责人 & 用户端与管理端基础工程；Element Plus 重写；响应式布局、导航和交互反馈；前端联调 \\
林洪宽 & \term{lin-hongkuan} & AI 与匹配负责人 & AI 服务、DashScope 接入、规则降级、图文匹配；匹配认领联调及部署协助 \\
郭焜华 & \term{gkh-gbh} & 后端 B 负责人 & 认领、交接、通知、积分、举报和管理模块；对应迁移、接口与测试 \\
程文韬 & \term{sptlayzner} & 文档与过程协作 & 与朱梓涵共同完成需求、设计、测试、运维及答辩材料；参与用例整理、验收检查和文档复核；期末报告提交标记二人共同完成 \\
\bottomrule
\end{longtable}

\subsection{职责协作关系}

图 \ref{fig:team-map} 表示主责与协作关系。项目负责人统一需求和接口基线，四个角色分别承担前端、后端 B、AI 匹配与文档过程，并在关键流程上共同联调。

\begin{figure}[H]
\centering
\resizebox{0.98\textwidth}{!}{%
\begin{tikzpicture}[x=1cm,y=1cm]
  \node[component,fill=sysugreen!14] (lead) at (0,2.2) {朱梓涵\\统筹、后端 A、集成};
  \node[component] (front) at (-6,0) {周星辰\\双前端};
  \node[component] (backb) at (-2,0) {郭焜华\\后端 B};
  \node[component] (ai) at (2,0) {林洪宽\\AI 与匹配};
  \node[component] (docs) at (6,0) {程文韬\\文档与过程};
  \draw[flow] (lead) -- (front);
  \draw[flow] (lead) -- (backb);
  \draw[flow] (lead) -- (ai);
  \draw[flow] (lead) -- (docs);
\end{tikzpicture}
}
\caption{团队主责与协作关系}
\label{fig:team-map}
\end{figure}

\clearpage
\section{阶段工作分配}

\zihao{-5}
\begin{longtable}{L{2.2cm}L{3.0cm}L{4.4cm}L{4.7cm}}
\caption{软件生命周期阶段分工}\label{tab:stages}\\
\toprule
阶段 & 主责成员 & 协作成员 & 主要产出 \\
\midrule
\endfirsthead
\toprule
阶段 & 主责成员 & 协作成员 & 主要产出 \\
\midrule
\endhead
需求与计划 & 朱梓涵、程文韬 & 全体成员 & 需求基线、优先级、WBS、进度与验收标准 \\
系统设计 & 朱梓涵、程文韬 & 林洪宽、郭焜华、周星辰 & 系统边界、数据模型、状态机与接口方案 \\
前端实现 & 周星辰 & 朱梓涵、郭焜华 & 用户端、管理端、共享类型、响应式界面与联调 \\
后端 A & 朱梓涵 & 郭焜华 & 用户认证、物品、上传、搜索与对象存储 \\
后端 B & 郭焜华 & 朱梓涵 & 认领交接、通知、积分、举报与后台管理 \\
AI 与匹配 & 林洪宽 & 朱梓涵、郭焜华 & AI 服务、规则增强、匹配评分与后端调用 \\
测试与修复 & 朱梓涵、程文韬 & 各模块负责人 & 测试计划、核心用例、质量修复与回归清单 \\
部署与交付 & 朱梓涵、林洪宽 & 程文韬、周星辰 & Compose、环境模板、CI、报告与演示材料 \\
\bottomrule
\end{longtable}
\normalsize

\section{成员具体贡献}

\subsection{朱梓涵}

\begin{itemize}
  \item 维护仓库规范、需求/架构基线、分支集成和交付节奏，组织全量质量检查与修复；
  \item 完成后端 A 的异步 FastAPI 基础设施、用户与认证、失物/招领、上传和对象存储等模块；
  \item 推进 durable job、AI 调用降级、敏感资源访问、管理统计、迁移和跨模块状态闭环；
  \item 建立并修复 GitHub Actions 门禁，覆盖后端、AI、双前端与 Docker smoke；
  \item 与程文韬共同编写需求、设计、测试、运维和课程交付文档，负责最终集成，并在期末报告提交中标记二人共同完成。
\end{itemize}

代表性提交主题包括 \term{feat(FR-01,FR-02)}、\term{fix(ci)}、\term{fix: 修复 P0/P1 交互与安全闭环} 和核心功能迭代。

\subsection{周星辰}

\begin{itemize}
  \item 初始化用户端和管理端 Vue 3/Vite/Pinia 工程，形成可独立构建的双前端；
  \item 使用 Element Plus 重写页面和交互，统一布局、状态反馈、表单体验和视觉规范；
  \item 改进移动端导航、返回流程、列表终态隐藏及页面反馈；
  \item 配合后端完成认证、物品、认领和管理接口联调，并补充前端工具函数测试。
\end{itemize}

\subsection{林洪宽}

\begin{itemize}
  \item 建立独立 AI FastAPI 服务及其配置、Schema、测试和 CI 基线；
  \item 接入 DashScope，同时提供可重复测试的规则型分类、敏感检测和匹配降级；
  \item 参与匹配认领闭环、主后端 AI 调用、匹配通知和异常降级联调；
  \item 协助部署环境、AI 服务镜像及相关架构文档。
\end{itemize}

\subsection{郭焜华}

\begin{itemize}
  \item 实现认领申请、问答、凭证、审核、申诉和双方交接状态流程；
  \item 实现通知、信誉积分、举报/后台管理和操作记录等后端 B 模块；
  \item 编写对应 ORM、Repository、Service、Router、Schema、迁移和单元测试；
  \item 与后端 A、前端和 AI 模块共同处理匹配到认领、交接到积分的业务闭环。
\end{itemize}

\subsection{程文韬}

\begin{itemize}
  \item 与朱梓涵共同梳理课程要求、需求基线、角色边界、功能优先级和验收标准；
  \item 参与系统模型、架构决策、测试计划、核心用例、缺陷清单和运维方案的撰写与复核；
  \item 将各模块负责人的实现信息整理为面向课程评审的文档和答辩材料；
  \item 参与最终验收清单检查，统一功能、测试和交付材料的完成口径。
\end{itemize}

课程文档由朱梓涵与程文韬共同完成。七份期末报告由朱梓涵负责最终集成，Git 提交使用 \term{Co-authored-by} trailer 标记程文韬为共同作者，使文档共创在版本记录中清晰呈现。

\section{协作过程与提交方式}

\subsection{文档共同提交}

代码提交由各模块负责人按职责完成；需求、设计、测试、运维和团队报告由朱梓涵与程文韬共同编写。最终报告提交同时记录 author 与共同作者 trailer，既保持统一集成，也完整体现文档贡献。

\subsection{协作实践}

\begin{enumerate}
  \item \textbf{任务分解：}WBS 将用户前端、管理前端、后端 A、后端 B、AI、测试与部署拆为可独立验收的工作包；
  \item \textbf{契约先行：}接口响应、枚举、状态机、匹配和积分规则先记录在 \path{docs/}，再由各模块实现；
  \item \textbf{分支与 PR：}功能分支合入 \term{develop}，稳定版本进入 \term{main}，通过 PR 暴露变更范围和讨论入口；
  \item \textbf{持续集成：}每次 push/PR 触发后端、AI、双前端和容器构建门禁，最近基线共 292 项测试通过；
  \item \textbf{集中集成：}项目负责人处理跨模块冲突、最终修复和文档提交；模块成员负责解释与验证自己的业务领域。
\end{enumerate}

\section{阶段成果与责任追踪}

\begin{longtable}{L{3.4cm}L{3.2cm}L{4.0cm}L{3.5cm}}
\caption{关键交付物责任追踪}\label{tab:deliverables}\\
\toprule
交付物 & 主责 & 评审/协作 & 验证方式 \\
\midrule
\endfirsthead
\toprule
交付物 & 主责 & 评审/协作 & 验证方式 \\
\midrule
\endhead
用户与认证、物品 API & 朱梓涵 & 郭焜华、周星辰 & 后端单测、前端构建、API 文档 \\
认领、交接、积分、通知 & 郭焜华 & 朱梓涵、周星辰 & 状态机测试、事务测试、页面联调 \\
用户端与管理端 & 周星辰 & 朱梓涵、郭焜华 & 32 项前端测试、双端生产构建 \\
AI 分类/敏感/匹配 & 林洪宽 & 朱梓涵 & 60 项 CI 测试、规则降级测试 \\
项目文档与报告 & 朱梓涵、程文韬 & 全体模块负责人 & 文档交叉检查、报告 PDF \\
CI 与部署环境 & 朱梓涵、林洪宽 & 程文韬 & GitHub Actions 五个 job、Compose smoke \\
质量回归与优化 & 朱梓涵 & 全体成员 & 检查清单、回归用例、绿色 CI \\
\bottomrule
\end{longtable}

\section{团队过程成效}

\subsection{做得较好的方面}

\begin{itemize}
  \item 模块边界与责任清晰，五个成员均有明确的主责领域；
  \item 统一技术栈和共享文档降低了前后端、AI 服务联调成本；
  \item 通过全量质量检查和持续回归，把核心流程推进为完整可演示闭环；
  \item 文档共创与模块复核同步进行，使七份课程报告与项目设计保持一致。
\end{itemize}

\subsection{协作经验与后续提升}

\begin{itemize}
  \item 延续 \term{Co-authored-by} 与独立 PR 方式，让代码、文档和评审贡献持续可追踪；
  \item 进一步缩小 PR 范围并强化交叉评审，使模块集成节奏更加均衡；
  \item 将完成定义、联调窗口和验收清单前移到每个迭代，保持稳定交付节奏；
  \item 继续扩展数据库、对象存储和浏览器旅程测试，丰富自动化测试层次；
  \item 使用 Issue 看板持续管理优化任务、负责人和完成时间。
\end{itemize}

\section{总结}

五名成员围绕同一业务闭环形成了前端、后端、AI、工程化和文档的互补分工。朱梓涵承担统筹、后端 A 和最终集成；周星辰负责双前端；林洪宽负责 AI 与匹配；郭焜华负责后端 B；程文韬与朱梓涵共同承担文档和过程材料。团队完成了项目代码、自动化验证、部署配置和七项课程报告，实现了可演示、可测试、可部署的课程项目目标。

\appendix
\section{项目资料索引}

\begin{itemize}
  \item 团队角色与 WBS：\path{docs/project-management/wbs-and-roles.md}
  \item 开发过程：\path{docs/project-management/development-plan.md}
  \item 需求基线：\path{docs/project-management/requirements-baseline.md}
  \item 协作规范：\path{CONTRIBUTING.md}、\path{AGENTS.md}、\path{CODEBUDDY.md}
  \item CI：\path{.github/workflows/ci.yml}
  \item 质量检查记录：\path{docs/testing/full-code-audit-2026-07-11.md}
  \item 版本记录：\verb|git shortlog -sne --all|、\verb|git log --all --stat|
\end{itemize}
